\newpage
% Оформляем заголовок приложения согласно СТП БГУИР
\begin{flushright}
    {\bfseries ПРИЛОЖЕНИЕ И} \\
    (обязательное)
\end{flushright}

\begin{center}
    {\bfseries РУКОВОДСТВО ОПЕРАТОРА} \\
    \vspace{0.2cm}
    БГУИР ДП 1-40 01 01 001 34 01
\end{center}

\addcontentsline{toc}{section}{Приложение И. Руководство оператора}

\section*{И.1 Назначение программы}

Программа «Система мониторинга параметров окружающей среды» предназначена для визуализации текущих и архивных данных о температуре и влажности, полученных от беспроводных узлов мониторинга. Пользовательский интерфейс реализован в виде Web-панели (Dashboard).

\section*{И.2 Условия выполнения программы}

Для корректной работы Web-интерфейса оператора необходимы:
\begin{itemize}
    \item персональный компьютер или планшет с выходом в сеть Интернет;
    \item современный браузер (Google Chrome версии 90+, Mozilla Firefox версии 88+, Safari 14+);
    \item разрешение экрана не менее 1280x720 пикселей.
\end{itemize}

\section*{И.3 Выполнение программы}

\subsection*{И.3.1 Запуск программы}
Для начала работы необходимо открыть браузер и ввести в адресной строке URL-адрес сервера мониторинга (например, \texttt{http://192.168.1.10} или адрес облачного сервиса). 

\subsection*{И.3.2 Описание интерфейса}
После загрузки страницы открывается главное окно мониторинга. Основные элементы интерфейса:
\begin{enumerate}
    \item \textbf{Индикаторы текущих значений} --- отображают последние полученные данные (обновление каждые 30 секунд).
    \item \textbf{Графики истории (Charts)} --- позволяют анализировать изменение параметров за последние 24 часа.
    \item \textbf{Статус соединения} --- отображает уровень сигнала Wi-Fi (RSSI) и время последнего сеанса связи.
\end{enumerate}



\subsection*{И.3.3 Настройка уведомлений}
Для получения уведомлений о превышении пороговых значений температуры (выше 28°C) необходимо перейти в раздел «Settings» и указать Email или Telegram ID в поле «Alert Notifications».

\section*{И.4 Сообщения оператору}

В процессе работы могут возникать следующие системные сообщения:
\begin{itemize}
    \item \textit{«Connection Lost»} --- потеряна связь с MQTT-брокером. Проверьте настройки сети.
    \item \textit{«Sensor Error»} --- данные от датчика не поступают. Требуется технический осмотр узла.
    \item \textit{«Battery Low»} --- напряжение питания ниже 3.3 В. Требуется подзарядка аккумулятора.
\end{itemize}

\section*{И.5 Аварийные ситуации}

При зависании Web-интерфейса необходимо выполнить принудительную перезагрузку страницы (клавиша F5 или Ctrl+R). В случае отсутствия данных на графиках более 1 часа следует перезагрузить центральный шлюз (Broker).

\vfill
\begin{center}
    Минск 2026
\end{center}
\newpage